% DIAPOSITIVA 1: Contexto Histórico y Filosofía
\begin{frame}{Contexto Histórico y Filosofía}

    \begin{columns}

        \begin{column}{0.65\textwidth}

            \textbf{Frank Plumpton Ramsey (1903--1930)} fue un genio
            británico cuyas ideas sentaron las bases de esta rama de
            la matemática.

            \vspace{0.4cm}

            \begin{exampleblock}{Principio de Motzkin}
                \centering
                \textit{``El desorden completo es imposible''.}
            \end{exampleblock}

            \vspace{0.3cm}

            \begin{block}{Concepto Clave}
                Si un sistema es lo suficientemente grande, por más
                caótico que parezca, \textbf{siempre aparecerán
                patrones} o subconjuntos con propiedades específicas.
            \end{block}

        \end{column}

        \begin{column}{0.35\textwidth}
            % Puedes insertar aquí una imagen de Ramsey, p.ej.:
            % \includegraphics[width=\linewidth]{ramsey.jpg}
        \end{column}

    \end{columns}

\end{frame}


% DIAPOSITIVA 2: Definición Formal en Teoría de Grafos
\begin{frame}{Definición Formal}

    \begin{alertblock}{Número de Ramsey $R(r,s)$}
        Para cualquier par de enteros positivos $(r,s)$, existe un
        número mínimo $R(r,s)$ tal que, si $n \geq R(r,s)$ y se
        colorean las aristas de un grafo completo $K_n$ con
        \textbf{dos colores}, entonces necesariamente aparece:
        \begin{itemize}
            \item un subgrafo $K_r$ con todas las aristas del
                  \textbf{primer color}, o
            \item un subgrafo $K_s$ con todas las aristas del
                  \textbf{segundo color}.
        \end{itemize}
    \end{alertblock}

    \pause

    \begin{center}
        \textbf{Valores conocidos:} \quad
        $R(3,3)=6$ \quad $R(3,4)=9$ \quad $R(3,5)=14$ \quad
        $R(4,4)=18$
    \end{center}

\end{frame}


% DIAPOSITIVA 3: Los Tres Pilares y Conclusión
\begin{frame}{Método de Deducción: Los Tres Pilares}

    Para demostrar o entender un problema de Ramsey se consideran:

    \begin{enumerate}

        \item \textbf{Estructura:} Se trabaja con un conjunto grande
              (un ``universo'' de elementos).

        \item \textbf{Partición (Conteo):} Los elementos se dividen
              en categorías —colores, grupos, clases.

        \item \textbf{Inevitabilidad (Orden):} Al alcanzar un tamaño
              crítico, la densidad fuerza la aparición de un
              \textbf{patrón uniforme}.

    \end{enumerate}

    \pause
    \vspace{0.3cm}

    \begin{block}{Conclusión}
        La Teoría de Ramsey muestra que \textbf{el azar tiene
        límites}. No importa cuánto se mezclen los elementos; si el
        sistema crece lo suficiente, los patrones aparecen por
        \textit{pura necesidad matemática}.
    \end{block}

\end{frame}